% !TeX root = booklet-01.tex
\section{مقدمات لازم برای مثال پردازنده \lr{MIPS}}
\subsection{پردازنده، حافظه، ثبات و دستور}
\lr{Processor} دستورها را از حافظه می‌خواند و طبق معنای آن‌ها داده‌ها و
حالت داخلی را تغییر می‌دهد. \lr{Memory} داده را در مکان‌های دارای آدرس نگه می‌دارد؛
آدرس، شماره مکان است و با محتوای آن مکان فرق دارد.
\lr{Register} محل کوچک ذخیره درون پردازنده و \lr{Register File}
مجموعه‌ای از این ثبات‌ها با مدار انتخاب و خواندن/نوشتن است.
این مقدمه برای خواندن بخش \secnum{۱.۷} کتاب لازم است \cite{weste}.

\lr{Instruction} یک الگوی بیت است که عملیات و عملوندها را مشخص می‌کند.
\lr{Program Counter (PC)} آدرس لازم برای خواندن جریان دستورها را نگه می‌دارد؛
در این طراحی، هنگام چهار \lr{Fetch} بایتی، به آدرس بایت بعدی پیش می‌رود.
\lr{Instruction Register (IR)} خود دستور خوانده‌شده را ذخیره می‌کند.
پس \lr{PC} آدرس است و \lr{IR} محتوا؛ جابه‌جا گرفتن این دو، مسیر داده را نامفهوم می‌کند.

دستور \code{add \$1, \$2, \$3} یعنی محتوای ثبات‌های دوم و سوم خوانده و
جمع در ثبات اول نوشته شود؛ اعداد نام ثبات‌اند، نه مقادیر داده.
اگر ثبات دوم ۵ و سوم ۷ داشته باشد، نتیجه ۱۲ است، نه جمع عددهای ۲ و ۳.
حروف \lr{Assembly} برای انسان خوانا هستند؛ سخت‌افزار کد دودویی متناظر را دریافت می‌کند.
\lr{Assembler} این کدگذاری را انجام می‌دهد؛ \lr{Compiler} زبان سطح بالا را
به نمایش پایین‌تر برنامه ترجمه می‌کند.

\subsection{عرض داده با عرض دستور یکسان نیست}
در زیرمجموعه کتاب، هشت ثبات عمومی هشت‌بیتی با نام‌های \code{\$0} تا \code{\$7}
و یک \lr{PC} هشت‌بیتی داریم؛ \code{\$0} همواره صفر است.
دستورها ۳۲ بیت و رابط داده حافظه هشت‌بیتی است
\cite[بخش \secnum{۱.۷.۱}]{weste}.
بنابراین ذخیره یک دستور در \lr{IR}، به چهار بایت هشت‌بیتی نیاز دارد:
\begin{equation}\label{eq:mips-fetch-width}
N_{\mathrm{bytes}}=\frac{32\ \mathrm{bits}}{8\ \mathrm{bits/byte}}=4\ \mathrm{bytes}.
\end{equation}
این نسبت، تعداد انتقال را می‌دهد؛ فرض شده هر چرخه یک بایت در دسترس است.
در طراحی دارای انتظار حافظه، تعداد چرخه می‌تواند بیشتر شود.

یک بیت دو حالت دارد؛ هشت بیت بدون علامت از صفر تا \(2^8-1=255\)
را نمایش می‌دهند. برای آدرس بایتی هشت‌بیتی، ۲۵۶ آدرس متمایز داریم.
اگر همه این فضا فقط دستورهای چهار‌بایتی باشد، حداکثر ۶۴ دستور جا می‌گیرد؛
ولی حافظه این مثال میان دستور و داده مشترک است و این تعداد ظرفیت آزاد واقعی برنامه را تضمین نمی‌کند.

برای نمایش منفی، قرارداد \lr{Two's Complement} را لازم داریم.
در هشت بیت، الگوی همه‌یک هم مقدار بدون علامت ۲۵۵ و هم مقدار علامت‌دار منفی یک
را می‌تواند نمایش دهد؛ تفسیر را عملیات تعیین می‌کند.
روش ساخت منفی یک از مثبت یک این است که بیت‌ها را مکمل و سپس یک اضافه کنیم:
\begin{latin}
\begin{verbatim}
 00000001   +1
 11111110   bitwise complement
 11111111   add one: -1 in signed 8-bit form
\end{verbatim}
\end{latin}
جمع هشت‌بیتی با نگه داشتن هشت بیت پایین، حساب پیمانه ۲۵۶ است؛
مثلاً الگوی ۲۵۵ به‌علاوه یک، هشت بیت پایین صفر و یک نقلی بیرونی می‌دهد.
این بیان حساب بیت‌هاست؛ پیاده‌سازی کامل استثناهای معماری \lr{MIPS32} را برای این مثال فرض نمی‌کنیم.

\section{مجموعه دستورها و کدگذاری}
\subsection{معنای ده دستور پشتیبانی‌شده}
کتاب دستورهای \code{ADD}، \code{SUB}، \code{AND}، \code{OR}، \code{SLT}،
\code{ADDI}، \code{BEQ}، \code{J}، \code{LB} و \code{SB} را معرفی می‌کند
\cite[جدول \secnum{۱.۷}]{weste}.
در جدول \ref{tab:mips-isa}، \(R_a,R_b,R_d\) محتوای ثبات، \(M[a]\)
بایت حافظه در آدرس \(a\) و \(\mathrm{imm}\) مقدار ثابت داخل دستور است.
پیکان \(\leftarrow\) یعنی نوشتن مقدار جدید و معادله تساوی دو مقدار هم‌زمان نیست.

\begin{table}[ht]
\centering
\caption{معنای دستورهای زیرمجموعه آموزشی؛ برای تغییر \lr{PC}، جزئیات متن را نیز بخوانید.}
\label{tab:mips-isa}
\begin{tabularx}{\textwidth}{lXl}
\toprule
دستور & اثر اصلی & \code{op}\\
\midrule
\code{ADD}&\(R_d\leftarrow R_a+R_b\)&\code{000000}\\
\code{SUB}&\(R_d\leftarrow R_a-R_b\)&\code{000000}\\
\code{AND}&عمل \lr{AND} بیت‌به‌بیت دو ثبات&\code{000000}\\
\code{OR}&عمل \lr{OR} بیت‌به‌بیت دو ثبات&\code{000000}\\
\code{SLT}&نوشتن یک در مقصد اگر شرط کمتر بودن برقرار باشد، وگرنه صفر&\code{000000}\\
\code{ADDI}&جمع ثبات با ثابت \lr{Immediate} و نوشتن در ثبات مقصد&\code{001000}\\
\code{BEQ}&اگر دو ثبات برابر باشند، انتقال به آدرس نسبی شاخه&\code{000100}\\
\code{J}&انتقال به مقصد پرش کدگذاری‌شده&\code{000010}\\
\code{LB}&\(R_b\leftarrow M[R_a+\mathrm{imm}]\)؛ خواندن یک بایت&\code{100000}\\
\code{SB}&\(M[R_a+\mathrm{imm}]\leftarrow R_b\)؛ نوشتن یک بایت&\code{101000}\\
\bottomrule
\end{tabularx}
\end{table}
در \code{LB} و \code{SB}، دو جمع متفاوت نداریم؛ ابتدا آدرس موثر با جمع
پایه و جابه‌جایی ساخته و سپس حافظه با همان آدرس استفاده می‌شود.
\code{AND} و \code{OR} روی هر جفت بیت متناظر مستقل عمل می‌کنند.
تفسیر علامت‌دار و رفتار سرریز در مقایسه و حساب باید در مشخصات پیاده‌سازی دقیق روشن باشد؛
مثال‌های عددی مقایسه این جزوه در محدوده بدون ابهام سرریز انتخاب می‌شوند.

\subsection{قالب‌های \lr{R}، \lr{I} و \lr{J}}
در هر سه قالب، شش بیت بالایی \code{op} هستند. قالب \lr{R} پنج فیلد دیگر دارد:
\code{rs}، \code{rt}، \code{rd}، \code{shamt} و \code{funct}.
سه فیلد ثبات هرکدام پنج بیت‌اند؛ در زیرمجموعه هشت‌ثباتی، فقط نام‌های مجاز
صفر تا هفت استفاده می‌شوند، هرچند محل کدگذاری پنج‌بیتی حفظ شده است.
فیلد پنج‌بیتی \code{shamt} در دستورهای \lr{R} مورد بحث صفر است.
قالب \lr{I} دو فیلد ثبات و ثابت ۱۶بیتی و قالب \lr{J} مقصد ۲۶بیتی دارد
\cite[شکل \secnum{۱.۴۹}]{weste}.

شمارش طول فیلدها، یک کنترل ساده بر کدگذاری است:
\begin{align}
\mathrm{R}:&\quad 6+5+5+5+5+6=32,\notag\\
\mathrm{I}:&\quad 6+5+5+16=32,\notag\\
\mathrm{J}:&\quad 6+26=32.
\label{eq:mips-fields}
\end{align}
دستورهای \lr{R} این مثال همگی \code{op=000000} دارند و \code{funct}
عملیات را جدا می‌کند؛ مثلاً \code{100000} برای جمع و \code{100010} برای تفریق.
قالب ۳۲بیتی را نباید به دلیل داده هشت‌بیتی به قالب کوتاه‌تر تبدیل کرد.

\subsection{مثال خواندن بیت‌ها و \lr{Hexadecimal}}
در مبنای شانزده، هر رقم چهار بیت را نمایش می‌دهد؛ پس ۳۲ بیت با هشت رقم
نمایش داده می‌شود. حروف \code{a} تا \code{f} به‌ترتیب مقادیر ده تا پانزده‌اند.
مثلاً \code{1111} برابر \code{f} است؛ این حروف بخشی از عددند، نه نام متغیر.

دستور \code{addi \$3, \$0, 8} به‌ترتیب فیلدهای زیر را دارد:
\begin{latin}
\begin{verbatim}
op     rs    rt    immediate
001000 00000 00011 0000000000001000
0010 0000 0000 0011 0000 0000 0000 1000
   2    0    0    3    0    0    0    8
\end{verbatim}
\end{latin}
پس کد \lr{Hex} برابر \code{20030008} است. مقدار ۸ در ثبات سوم نوشته می‌شود،
چون ثبات صفر مقدار صفر می‌دهد. در \code{addi \$5, \$0, -1}، ثابت ۱۶بیتی
\code{ffff} است؛ هشت بیت پایین آن در مسیر داده هشت‌بیتی، الگوی \code{ff} را می‌دهد.
این کوچک شدن مسیر داده، فیلد ذخیره‌شده ۱۶بیتی داخل دستور را تغییر نمی‌دهد.

\subsection{آدرس شاخه و پرش را با واحد بایت بخوانیم}
دستورها چهار بایت‌اند؛ مقصد نسبی شاخه در واحد کلمه دستور کدگذاری می‌شود.
پس باید جابه‌جایی در چهار ضرب شود؛ این کار با دو بیت شیفت چپ متناظر است.
اگر آدرس شروع دستور \(P\) باشد، پس از \lr{Fetch} چهار بایت، \lr{PC} به
\(P+4\) رسیده است؛ در صورت برابری دو ثبات، مقصد شاخه در این مثال برابر است با:
\begin{equation}\label{eq:mips-branch}
PC_{\mathrm{target}}=(P+4+4\,\mathrm{imm})\bmod256.
\end{equation}
نماد \(\bmod256\) یعنی نگه داشتن باقیمانده در بازه صفر تا ۲۵۵،
متناسب با \lr{PC} هشت‌بیتی. برای \(\mathrm{imm}\) منفی،
تفسیر علامت‌دار جابه‌جایی لازم است؛ محاسبه پیمانه‌ای نتیجه بیت‌های پایین را مشخص می‌کند.
استفاده از \(P\) به‌جای \(P+4\) یک دستور خطای مقصد ایجاد می‌کند.

در برنامه کتاب، شاخه در آدرس ۱۲ و برچسب پایان در آدرس ۳۲ است؛
پس مقدار جابه‌جایی دستور شاخه چنین محاسبه می‌شود.
\[
 \mathrm{imm}=\frac{32-16}{4}=4.
\]
برای پرش به آدرس ۱۲، فیلد مقصد \(12/4=3\) است؛
سخت‌افزار با شیفت چپ دو بیت، آن را به آدرس بایتی برمی‌گرداند.
قاعده ترکیب بیت‌های بالایی مقصد در معماری کامل ۳۲بیتی را نباید بدون توجه
به \lr{PC} هشت‌بیتی این مدل به شکل مثال تحمیل کرد
\cite[جدول \secnum{۱.۷}، شکل‌های \secnum{۱.۵۲} و \secnum{۱.۵۳}]{weste}.

\section{مثال \lr{Fibonacci}: از برنامه تا کد ماشین}
\subsection{مقادیر اولیه و ترتیب به‌روزرسانی}
مثال \secnum{۱.۴} کتاب تابع بازگشتی را با
\(f_{-1}=-1\)، \(f_0=1\) و \(f_n=f_{n-1}+f_{n-2}\) برای \(n>0\)
تعریف می‌کند. این مقادیر اولیه با قرارداد رایج \(F_0=0,F_1=1\) یکسان نیستند؛
پس نام \lr{Fibonacci} به‌تنهایی پاسخ عددی را تعیین نمی‌کند
\cite[شکل \secnum{۱.۵۰}]{weste}.
در برنامه، \code{f1} ابتدا یک و \code{f2} منفی یک است و حلقه هشت بار اجرا می‌شود:
\begin{latin}
\begin{verbatim}
f1 = f1 + f2;
f2 = f1 - f2;
n  = n - 1;
\end{verbatim}
\end{latin}
این سه خط نرم‌افزاری به‌ترتیب اجرا می‌شوند. اگر مقادیر قدیمی را \(a,b\)
بنامیم، بعد از خط اول \(f1=a+b\) و بعد از خط دوم
\(f2=(a+b)-b=a\) می‌شود. بنابراین زوج به \((a+b,a)\) تبدیل می‌شود.
اگر در خط دوم به اشتباه از \code{f1} قدیمی استفاده کنیم، الگوریتم دیگری خواهیم داشت.

جدول \ref{tab:mips-fib} نتیجه اجرای مستقل همین ترتیب را نشان می‌دهد؛
در این ورودی، مقادیر نتیجه از محدوده علامت‌دار هشت‌بیتی خارج نمی‌شوند.
\begin{table}[ht]
\centering
\caption{ردگیری آموزشی حلقه با مقادیر اولیه دقیق کتاب؛ ستون‌ها مقدار پس از تکرار را نشان می‌دهند.}
\label{tab:mips-fib}
\begin{tabular}{cccc}
\toprule
تعداد تکرار اجراشده & شمارنده باقی‌مانده & \code{f1} & \code{f2}\\
\midrule
۰&۸&۱&\(-1\)\\
۱&۷&۰&۱\\
۲&۶&۱&۰\\
۳&۵&۱&۱\\
۴&۴&۲&۱\\
۵&۳&۳&۲\\
۶&۲&۵&۳\\
۷&۱&۸&۵\\
۸&۰&۱۳&۸\\
\bottomrule
\end{tabular}
\end{table}
پس مقدار ذخیره‌شده نهایی ۱۳ است. برنامه صرفاً نام دنباله را به کار نبرده،
بلکه قرارداد دقیق اولیه و عملیات ترتیب‌دار دارد.

\subsection{خواندن برنامه \lr{Assembly}}
در شکل \secnum{۱.۵۱} کتاب، ثبات سوم شمارنده، چهارم \code{f1} و پنجم
\code{f2} است. برچسب‌ها نشانی‌اند و خودشان دستور اجرایی نیستند.
برنامه زیر همان دستورها را با توضیح انگلیسی کوتاه بازنویسی می‌کند:
\begin{latin}
\small
\begin{verbatim}
fib:  addi $3, $0, 8     # n = 8
      addi $4, $0, 1     # f1 = 1
      addi $5, $0, -1    # f2 = -1
loop: beq  $3, $0, end   # exit when n == 0
      add  $4, $4, $5    # update f1 first
      sub  $5, $4, $5    # use the new f1
      addi $3, $3, -1    # decrement n
      j    loop
end:  sb   $4, 255($0)   # write result to byte address 255
\end{verbatim}
\end{latin}
وقتی شمارنده صفر شود، \code{beq} به \code{end} می‌رود و یک بایت نتیجه را
در آدرس ۲۵۵ می‌نویسد. صفر بودن \code{\$0} آدرس موثر \(0+255\) را می‌دهد.
این برنامه دستور \code{HALT} ندارد؛ مشاهده نوشتن نتیجه، به معنی توقف ذاتی پردازنده نیست.

\subsection{تطبیق همه کدهای ماشین مثال}
جدول \ref{tab:mips-machine} کدهای شکل \secnum{۱.۵۲} کتاب را با آدرس بایتی
دستورها کنار هم می‌گذارد. فرض این جدول، شروع برنامه از آدرس صفر است.
\begin{table}[ht]
\centering
\caption{کدهای ماشین برنامه کتاب؛ رشته‌های کد و شناسه‌ها با ارقام \lr{Latin} اصلی حفظ شده‌اند.}
\label{tab:mips-machine}
\begin{tabular}{cll}
\toprule
آدرس بایتی & دستور & کد \lr{Hex}\\
\midrule
۰&\code{addi \$3, \$0, 8}&\code{20030008}\\
۴&\code{addi \$4, \$0, 1}&\code{20040001}\\
۸&\code{addi \$5, \$0, -1}&\code{2005ffff}\\
۱۲&\code{beq \$3, \$0, end}&\code{10600004}\\
۱۶&\code{add \$4, \$4, \$5}&\code{00852020}\\
۲۰&\code{sub \$5, \$4, \$5}&\code{00852822}\\
۲۴&\code{addi \$3, \$3, -1}&\code{2063ffff}\\
۲۸&\code{j loop}&\code{08000003}\\
۳۲&\code{sb \$4, 255(\$0)}&\code{a00400ff}\\
\bottomrule
\end{tabular}
\end{table}
برای کنترل مستقل، می‌توان هر کد را به فیلدها شکست؛ مثلاً \code{00852020}
ثبات‌های منبع چهار و پنج، مقصد چهار و \code{funct=100000} را می‌دهد.
برای شاخه نیز آدرس ۱۲، \lr{PC} بعد از \lr{Fetch} برابر ۱۶ و جابه‌جایی چهار کلمه
با مقصد ۳۲ سازگار است.
